% Appendix B -- the correctness ladder, moved out of section 6 (2026-09-26) for
% the page limit. A reviewer who doubts the mechanism should read this; one who
% doubts the economics should not need to.

\section{Correctness ladder}
\label{app:correct}

Placement moves bytes rather than approximating them, so these gates are
hygiene rather than a contribution. Numerically we verify the host pool's round
trip, the two-tier attention against a single-pass reference over the same KV
(576 mixed decode layer-steps, $\max|\mathrm{diff}| = 0$), and the prefill
handover across block boundaries and after the producer's memory has been
recycled. End to end we compare generated tokens against stock vLLM.

One methodological note, because it traps anyone reproducing this.
\texttt{bf16} attention reductions are not associative, so their result depends
on batch composition: stock vLLM produces different text for the same prompt at
concurrency 1 and at 6, reproducibly at each. Any system that changes batching
--- as a memory tier does --- will diverge on such prompts without being wrong.
We therefore compare serially, batch of one on both sides, and under load only
on prompts pre-screened as batching-stable. For calibration: on a serial greedy
comparison our host tier matched stock vLLM on 6 of 6 prompts in one round and
4 of 6 in another, while enabling \emph{vLLM's own prefix cache} --- changing
nothing else --- matched on 2 of 6.


\section{Evaluation setup in full}
\label{app:setup}

\paragraph{Exclusive access, for every absolute.} Each number we report as an
absolute is taken on an idle, exclusively allocated device. This is not
boilerplate: the same measurement on a contended node and on an idle one
differed by $9.4\times$ (Appendix~\ref{app:impl}). Ratios between arms run back
to back are robust to contention; absolutes are not.

\paragraph{Models.} Qwen2.5-32B fp16 as the primary model, because on a 96\,GB
device its weights leave $\Ckv \approx 26$\,GB, the regime \S\ref{sec:model}
identifies as favourable. \textbf{On the A100-80 we actually run, the same
weights leave a KV pool of $11.8$\,GB}, so $B_{\mathrm{cap}} \approx 22$ at 2048
context: capacity-bound, as the model requires, but at a smaller
$\HBM/\Ckv$ than the projections assume. Numbers here should be read as coming
from that substrate, not the projected one. Llama-3.1-8B
as a negative control, where the model predicts no benefit. We deliberately do
not evaluate a 70B-class model: its capacity-limited batch is about seven
requests, so a single request is 14\% and no result could separate an effect
from a rounding artifact.

\paragraph{Workloads.} Every number reported in this section comes from a
\emph{fixed-context, fixed-output, closed-loop} run: 2048 context and 64 output
tokens, with the offered load and resident cohort stated per experiment. That is
the right shape for isolating a rider's cost and the wrong shape for goodput,
which is why \S\ref{sec:eval:goodput} is a protocol rather than a result. That
protocol calls for ShareGPT and Alpaca under Poisson arrivals, via
\texttt{vllm bench serve}. Two properties are required rather than preferred.
Real traces are heavy-tailed in both prompt and generation length, which is the
variation a \emph{per-request} placement policy exists to exploit --- on a
workload of identical requests our placement rule assigns every request the
same score and degenerates to arrival order by construction, so measuring it
equal to arrival order there would say nothing about the policy. And open-loop
arrivals are needed because a closed loop of $N$ simultaneous requests is a
different experiment from offering $\lambda$ requests per second; it also
invites metrics that reward slowness (below). These are the workloads Pie
reports on~\cite{pie}, which makes our numbers comparable to theirs.

\paragraph{Metric.} The goodput comparison of \S\ref{sec:eval:goodput} is
defined on \textbf{SLO-attaining goodput}: requests per second whose
per-output-token latency met a target. That experiment is not yet run, so no
goodput number appears in this section; what we report here is step time,
end-to-end decode throughput, achieved concurrency and transferred bytes, each
named where it is used.

We do \emph{not} report mean concurrency. It averages over sampled steps, so an
arm that serves requests more slowly holds them longer and scores
\emph{higher}. It is not a goodput metric, and an earlier result of ours that
relied on it was retracted for exactly that reason.

\paragraph{Baselines.} Stock vLLM with all tiering disabled, which is also
Pie's baseline and makes the two sets of numbers commensurable. We verify separately that our fork with the feature off matches the unmodified
engine: across six interleaved runs on Qwen2.5-0.5B it is slower by $0.08\%$,
% src: paper/data/baseline_integrity_20260912.log @ 5377be7c3 less than the
spread within either arm. A baseline quietly slowed by our own additions would
inflate every ratio reported here, and it would do so invisibly. Uniform offload at a matched byte budget represents the placement
that Pie and DirectKV instantiate at their own granularities. Where a
comparison against a published figure is drawn, we state which engine version
and workload each side used rather than presenting ratios as directly
commensurable.

\subsection*{Swap semantics versus streaming}

\paragraph{Swap semantics versus streaming.} A design whose offloaded unit is
\emph{owned} by one side must write that unit back once per generated token; one
whose device-side copy is a \emph{cache} need not (\S\ref{sec:related}). We
measure the difference directly by writing each layer's slot back to the host
mirror with the policy otherwise unchanged. The emulation is exact in extent:
both arms read the same $602$\,MB, and the write-back arm moves
\textbf{$2.00\times$} the total link traffic, with token output verified
identical. We report this as a mechanism check rather than a performance
result: it is measured on Qwen2.5-0.5B at $512$ context, a configuration small
enough that the ratio isolates the write-back leg without the step-time effects
that dominate at the primary model's scale.
% src: paper/data/ablation_swapback_20260911.json @ 74a5d9c58
